\documentclass[11pt,a4paper]{article}
\usepackage[T1]{fontenc}
\usepackage[utf8]{inputenc}
\usepackage[french]{babel}
\usepackage{amsmath,amssymb,booktabs,longtable,array}
\usepackage[a4paper,margin=2.35cm]{geometry}
\usepackage[hidelinks]{hyperref}
\usepackage{microtype}
\setlength{\parindent}{0pt}
\setlength{\parskip}{0.5em}
\newcommand{\NW}{\mathrm{NW}}
\newcommand{\code}[1]{\texttt{#1}}

\title{Spécification technique du modèle M2\\
\large Sémantique, invariants et protocole de validation}
\author{Réplication computationnelle indépendante}
\date{3 juillet 2026}

\begin{document}
\maketitle

\begin{abstract}
Ce document fige la sémantique de l'implémentation de référence avant son
écriture. Le noyau sera autonome sous \code{m2\_codex/}, sans modification des
lignées WIP. Les décisions visent une dynamique déterministe conditionnellement
au seed, une comptabilité vérifiable phase par phase et des sorties permettant
de réfuter les hypothèses Boltzmann, Pareto, cohorte et SOC. Les ambiguïtés du
rapport ne sont pas dissimulées : service de plusieurs dettes, liquidation sans
créancier et ordre des faillites deviennent des conventions nommées et des
ablations.
\end{abstract}

\tableofcontents

\section{Objectifs et périmètre}

Le premier objectif est la fidélité aux sept phases M2. Le second est la
traçabilité : toute quantité publiée doit être reconstructible à partir de la
configuration, du seed, des séries et snapshots. Le troisième est la prudence
scientifique : le code ne doit contenir aucun mécanisme destiné implicitement à
améliorer un ajustement statistique.

À la demande de séparation avec un autre agent, tous les nouveaux fichiers sont
placés sous \code{m2\_codex/}. L'arborescence proposée est conservée relativement
à ce dossier :
\begin{verbatim}
anciens_modeles/m2_codex/
  src/m2/              noyau, mesures, analyses, figures
  experiments/m2/      baseline, validation, grille, ablations
  tests/                tests unitaires et de cohérence
  outputs/              exécutions régénérables, une par run
  reports/              sources LaTeX, PDF, figures et synthèses
\end{verbatim}
Le répertoire de travail isolé est la seule adaptation à l'arborescence demandée.

\section{Architecture et structures de données}

\subsection{Modules}

\begin{longtable}{p{0.23\textwidth}p{0.69\textwidth}}
\toprule
Module & Responsabilité unique \\
\midrule
\endhead
\code{config.py} & Dataclass immuable \code{M2Config}, validation des domaines,
sérialisation et variantes structurelles nommées. \\
\code{entities.py} & \code{Entity} : identifiant, capital, dette d'amorçage,
naissance, état vivant/défaut et flux du dernier pas. \\
\code{contracts.py} & \code{LoanContract} nominal et \code{ContractBook} indexé
par prêteur/emprunteur ; création, suppression, fractionnement et contrôles. \\
\code{market.py} & Taux marginaux, cible $w^*$ et rounds de marché ; aucun effet
de faillite. \\
\code{bankruptcy.py} & Détection, liquidation unitaire, cascade et événements de
récupération/destruction. \\
\code{simulation.py} & État global, RNG NumPy local, sept phases, orchestration
et vérification optionnelle des invariants. \\
\code{metrics.py} & Séries temporelles, snapshots individuels, mobilité,
concentration du crédit et diagnostics anti-cohorte. \\
\code{analysis.py} & Ajustements du corps, CSN, comparaison log-normale,
bootstrap par blocs et estimation des incréments. \\
\code{plots.py} & Figures à partir de fichiers de résultats, sans relancer la
simulation ni modifier les données. \\
\bottomrule
\end{longtable}

\subsection{Entité}

La structure minimale est :
\begin{verbatim}
Entity(id: int, w: float, d0: float, birth_step: int,
       alive: bool, defaulted: bool,
       extraction: float, interest_received: float,
       interest_paid: float)
\end{verbatim}
Les créances et dettes ne sont pas des caches modifiables dans l'entité : elles
sont calculées par le carnet ou maintenues dans des index contrôlés. Cela évite
la divergence entre contrat et bilan. L'âge au snapshot $t$ vaut
$t-\text{birth\_step}$. La valeur nette est une propriété dérivée
\[
\NW_i=w_i+C_i-D_i-d_{0,i}.
\]

\subsection{Contrat et carnet}

\begin{verbatim}
LoanContract(id: int, lender_id: int, borrower_id: int,
             principal: float, rate: float, creation_step: int,
             parent_id: int | None)
\end{verbatim}
Un contrat présent dans le carnet est actif ; aucune double notion
actif/inactif. Les index \code{by\_lender} et \code{by\_borrower} contiennent des
ensembles d'identifiants. Les identifiants ne sont jamais recyclés. Un transfert
de créance supprime le parent et crée des enfants dont \code{parent\_id} permet
l'audit.

\subsection{État de simulation}

\code{M2Simulation} contient la configuration, le pas courant, un générateur
\code{numpy.random.Generator}, le dictionnaire d'entités, le carnet, les compteurs
d'identifiants, les métriques et les événements. Aucune variable aléatoire
globale n'est utilisée. Un même couple configuration--seed doit produire des
fichiers de contenu identique, hors métadonnées temporelles explicitement
exclues.

\section{Sémantique exacte d'un pas}

\subsection{Conventions de temps et unités}

Un pas est une unité discrète indivisible : taux $r$, dépréciation $\delta$ et
écart-type $\sigma$ sont par pas. $\alpha=1$ fixe l'unité conjointe de capital et
de flux ; il reste configurable seulement pour tester l'homogénéité, sans être
balayé dans la baseline. Les valeurs de capital, principaux, intérêts et $d_0$
partagent la même unité nominale. Un taux n'a pas cette unité.

Le snapshot $t$ est pris après la cascade du pas $t$ ; les naissances du pas ont
alors un âge nul si le compteur n'est incrémenté qu'après écriture, convention à
tester. Les fenêtres $[a,b)$ incluent les snapshots tels que $a\leq t<b$.

\subsection{Phases}

\begin{enumerate}
  \item Tirer $n\sim\mathrm{Poisson}(\lambda)$ et créer $n$ entités à
  $(w_0,d_0)$.
  \item Pour chaque vivante, tirer
  $\eta\sim\mathcal N(-\sigma^2/2,\sigma^2)$ et poser
  $w\leftarrow w\exp\eta$. Aucun clipping silencieux : une valeur non finie
  arrête le run avec diagnostic.
  \item Calculer $g_i=\alpha\sqrt{w_i}$, enregistrer l'extraction et poser
  $w_i\leftarrow w_i+g_i$.
  \item Grouper les intérêts dus par emprunteur. Si le capital couvre la somme,
  payer chaque $rq$. Sinon, partager tout le capital disponible entre les
  créanciers au prorata des montants dus et marquer l'emprunteur en défaut. Ce
  choix simultané évite qu'un ordre de dictionnaire crée une priorité. Une
  variante \code{sequential\_contract\_id} sera disponible pour l'ablation.
  \item Poser $w_i\leftarrow(1-\delta)w_i$ pour les vivantes. Contrats et $d_0$
  sont inchangés.
  \item Former l'ensemble des vivantes non marquées en défaut. Effectuer
  $\lfloor N/k\rfloor$ rounds indépendants ; échantillon sans remise au sein
  d'un round, mais participation répétée possible entre rounds. En cas d'égalité,
  le plus petit identifiant départage le max/min pour la reproductibilité. Si
  $w_b=0$ ou $w_\ell=0$, la limite est définie comme absence de prêt. Transférer
  strictement $q>0$ et créer le contrat.
  \item Construire la file des défauts puis des $\NW<0$, triée par identifiant.
  Liquider une entité à la fois, rescanner les insolvabilités, et continuer
  jusqu'à stabilité.
\end{enumerate}

L'exclusion des entités en défaut du marché n'est pas explicitée dans le rapport,
mais elle évite de créer un prêt destiné à être annulé dans la même phase. Ce
choix est journalisé comme convention structurelle.

\subsection{Liquidation de référence}

Avant toute suppression, agréger les contrats où la faillie est emprunteuse et
calculer les poids de ses créanciers vivants. Pour chaque contrat où elle est
prêteuse : supprimer le contrat parent, exclure comme repreneur l'emprunteur de
ce contrat afin d'éviter une auto-créance, renormaliser les poids admissibles,
puis créer les fractions nominales. Sans repreneur admissible, annuler le
contrat : l'emprunteur est libéré et la créance est détruite.

Distribuer ensuite le capital résiduel aux créanciers vivants, selon les poids
renormalisés. En l'absence de créancier vivant, le résiduel est détruit et
comptabilisé. Enfin annuler les dettes de la faillie, éteindre son $d_0$, poser
$w=0$ et la marquer morte. Une entité morte ne peut apparaître dans aucun contrat.

Cette règle suit l'intention du rapport et la logique du WIP, mais l'ordre
séquentiel peut influer sur une vague. Une variante de liquidation simultanée
fera partie des tests de sensibilité, pas de la baseline.

\section{Configuration et paramètres}

\begin{longtable}{p{0.24\textwidth}p{0.16\textwidth}p{0.50\textwidth}}
\toprule
Champ & Baseline & Validation \\
\midrule
\endhead
\code{alpha} & 1 & strictement positif ; convention \\
\code{delta} & 0,05 & $0\leq\delta<1$ \\
\code{k} & 6 & entier $\geq2$ \\
\code{sigma} & 0,25 & $\geq0$ \\
\code{birth\_rate} & 10 & $\geq0$ \\
\code{w0}, \code{d0} & 30, 28 & $w_0\geq d_0\geq0$ ; baseline exige $>$ \\
\code{seed} & explicite & entier non négatif \\
\code{steps} & 2000 & entier non négatif \\
\code{snapshot\_interval} & 50 & entier positif \\
\code{interest\_service} & prorata & variante nommée seulement \\
\code{bankruptcy\_transfer} & prorata & variantes annulation/destruction \\
\code{check\_invariants} & vrai & technique, sans effet dynamique \\
\bottomrule
\end{longtable}

La fréquence des snapshots ne doit pas affecter le RNG ni la dynamique. Les
tolérances flottantes sont centralisées et réservées aux assertions ; elles ne
suppriment pas de contrats positifs et ne changent pas le critère scientifique
$\NW<0$.

\section{Invariants et audits de conservation}

Les invariants sont vérifiés après chaque phase en mode test et périodiquement
dans les campagnes :
\begin{enumerate}
  \item toute entité vivante a un $w$ fini et $w\geq0$ ;
  \item tout contrat actif a $q>0$, $r\geq0$, deux extrémités vivantes et
  distinctes ;
  \item chaque contrat apparaît exactement une fois dans chacun des deux index ;
  \item $\sum_i C_i=\sum_i D_i=\sum_e q_e$ à l'erreur d'arrondi près ;
  \item la création d'un prêt conserve $\sum_iw_i$ et chaque $\NW_i$ concernée ;
  \item le service des intérêts conserve $\sum_iw_i$, y compris en paiement
  partiel ; l'impayé n'est ni capitalisé ni créé ailleurs ;
  \item les principaux, taux et $d_0$ ne changent pas à la dépréciation ;
  \item toute entité avec $\NW<0$ ou en défaut est retirée par la cascade ;
  \item à la fin d'une cascade, aucune vivante n'est défaillante ou insolvable ;
  \item après liquidation, aucun contrat orphelin ne subsiste ;
  \item il n'existe aucune borne, capacité ou suppression dépendant directement
  de $N$ ; $\lfloor N/k\rfloor$ ne règle que l'intensité du marché.
\end{enumerate}

Un audit de phase écrit les variations observées de capital réel. Naissances,
choc, extraction, dépréciation et destruction en liquidation sont les seules
sources/puits ; principal et intérêts doivent avoir une variation globale nulle.

\section{Tests avant simulations longues}

La suite utilisera le module standard \code{unittest}, disponible sans dépendance,
avec commande unique
\code{PYTHONPATH=src python -m unittest discover -s tests -v}. Le venv actuel ne
contient pas \code{pytest}; l'ajouter n'est pas nécessaire à la validité des tests.

\begin{longtable}{p{0.31\textwidth}p{0.60\textwidth}}
\toprule
Famille & Cas exigés \\
\midrule
\endhead
Entités et bilans & naissance, âge, $\NW$, $w=0$, dette d'amorçage. \\
Contrats & création/suppression, index, $q=0$ refusé, principal positif,
absence d'orphelin. \\
Conservation & transfert du principal, intérêt complet, intérêt partiel prorata,
dépréciation nominale nulle. \\
Défaut & marquage, faillite simple avec/sans créancier, récupération, extinction
de $d_0$. \\
Cascade & perte de créance rendant un prêteur insolvable, itération jusqu'à
stabilité, transfert fractionné sans auto-contrat. \\
Marché & $N=0$, $N<k$, égalités, $w=0$, $q=0$, volume et taux connus. \\
Reproductibilité & mêmes seed/configuration identiques ; seeds distincts non
identiques ; snapshots sans effet sur trajectoire. \\
Numérique & aucun capital négatif, aucune valeur non finie, longue séquence
smoke-test avec invariants. \\
Statistique & ajustement exponentiel synthétique, Pareto synthétique, rejet ou
prudence sur échantillon trop petit, fenêtres non mélangées. \\
\bottomrule
\end{longtable}

Aucune simulation de 2000 pas ne sera lancée tant que ces tests de base ne sont
pas tous verts.

\section{Risques numériques et performance}

Les principaux risques sont : overflow de $\exp(\eta)$ ; perte de précision dans
les longues sommes de contrats ; explosion du nombre de petits contrats après
fractionnements ; comparaison de $\NW$ près de zéro ; et complexité du rescannage
des cascades. Aucune troncature de richesse ni fusion de contrats ne sera ajoutée
à la baseline, car elle changerait la dynamique. Les sommes d'audit utiliseront
\code{math.fsum}. Les non-finis provoquent un échec explicite. Les index limitent
les opérations de carnet au voisinage de l'entité ; le rescannage $O(N)$ après
chaque vague est accepté d'abord pour sa clarté, puis profilé.

\section{Sorties, logs, snapshots et figures}

Chaque run a un identifiant fondé sur scénario, seed et horodatage humain, dans
un dossier contenant :
\begin{itemize}
  \item \code{run.json} : version de schéma, configuration complète, seed,
  commande, versions Python/NumPy/SciPy, statut et résumé ;
  \item \code{timeseries.csv} : une ligne par pas, avec population, naissances,
  décès, taille maximale de cascade, capital total, volumes et nombre de prêts,
  contrats, intérêts, destructions et concentration du crédit ;
  \item \code{snapshots/snapshot\_TTTTTT.csv} :
  \code{step,entity\_id,birth\_step,age,w,claims,debts,d0,nw,extraction,
  interest\_received,interest\_paid,income\_gross,income\_net} ;
  \item \code{events.jsonl} : événements optionnels de prêt, défaut et faillite,
  avec identifiants et montants ;
  \item \code{analysis.json} et \code{tables/*.csv} : estimations, domaines,
  effectifs, incertitudes et verdicts ;
  \item \code{figures/*.png} et \code{*.pdf} : mêmes données, noms stables,
  légende incluant scénario, seed et fenêtre.
\end{itemize}

Les figures attendues couvrent macro-séries, distributions et CCDF log--log,
QQ exponentiels, âge par décile, mobilité du top, cascades, concentration du
crédit et cartes de grille. Une figure ne doit jamais agréger des seeds sans
montrer la dispersion inter-seed.

\section{Protocole de validation}

\subsection{Baseline et stationnarité}

Lancer d'abord un smoke-test court, puis la baseline
$(\alpha,\delta,k,\sigma,\lambda,w_0,d_0)=(1,0{,}05,6,0{,}25,10,30,28)$ sur
2000 pas. Répliquer au moins trois seeds si le coût est acceptable. Diagnostiquer
la pente de population et de capital sur $[1000,2000)$ avec incertitude HAC ou
par blocs, le bilan naissances--décès, la survie censurée et la mortalité du top.
Un plateau visuel seul ne suffit pas.

\subsection{Corps}

Pour $w$, $\NW>0$, revenu brut et net, définir le corps par plusieurs quantiles
(90, 95 et 97,5\,\%), sans choisir après inspection. Ajuster exponentielle à
origine fixée, gamma, log-normale et Fisk par maximum de vraisemblance ; rapporter
AIC, BIC, effectif, paramètres et sensibilité au domaine. Produire QQ-plot et
ratio médiane/moyenne. Sur les survivantes présentes à deux snapshots successifs,
estimer dans le corps initial
\[
 \widehat A_0=-\mathbb E[\Delta\NW],\qquad
 \widehat B_0=\tfrac12\operatorname{Var}(\Delta\NW),
\]
et comparer $T_{\rm exp}$ à $\widehat B_0/\widehat A_0$ seulement si
$\widehat A_0>0$. Rapporter le biais de sélection dû aux morts exclues.

\subsection{Queue}

Pour chaque seed séparément, pooler uniquement des snapshots espacés de 50 pas
dans $[500,1000)$, $[1000,1500)$ et $[1500,2000)$. Estimer $x_{\min}$ et
l'exposant par MLE--KS de type CSN, puis comparer Pareto et log-normale sur le
même domaine. Répéter à $x_{\min}$ commun, médiane des seuils des trois fenêtres.
Le bootstrap rééchantillonne des blocs temporels/snapshots, non des lignes
individuelles supposées indépendantes. Une queue de moins de 100 observations est
déclarée fragile même si un estimateur existe.

\subsection{Anti-cohorte}

Mesurer corrélations Pearson et Spearman âge--$\log\NW$ sur $\NW>0$, âges par
décile, formes dans les tranches d'âge $[50,150)$ et $[150,400)$, part récente du
top décile, mortalité du top et Jaccard/turnover du top entre fenêtres. Ajouter
une régression ou stratification du risque de décès par âge et décile : une faible
corrélation marginale ne suffit pas.

\subsection{SOC, grilles et ablations}

Balayer $k\in\{2,3,4,6\}$, $\sigma\in\{0{,}15,0{,}25,0{,}35\}$ et
$\varepsilon_0/w_0\in\{0{,}05,0{,}25,0{,}5\}$, sans pooling entre cases. Mesurer
distribution des cascades, concentration de crédit, corrélation concentration
pré-cascade--taille et exposant de richesse. Les ablations minimales sont : sans
crédit, sans choc, sans transfert de créances, service séquentiel et liquidation
simultanée. Pour parler de SOC, rechercher aussi sensibilité à la taille via
$\lambda$, coupures de cascade et séparation des échelles ; une simple droite
log--log n'est pas une preuve.

\section{Résultats de spécification, limites et décisions ouvertes}

Le résultat de cette phase est une sémantique testable et non encore un résultat
scientifique. Les choix prorata et ordre par identifiant rendent la référence
reproductible, mais peuvent affecter les cascades. Le format CSV privilégie
l'audit au volume ; il pourra devenir compressé sans changer le schéma. La
comparaison Pareto--log-normale existante du dépôt n'utilise pas une log-normale
correctement tronquée au même $x_{\min}$ ; elle devra être corrigée localement
avant toute conclusion. L'absence de \code{pytest} dans le venv est contournée
par \code{unittest}, non par installation implicite d'une dépendance.

Les prochaines étapes sont : implémenter d'abord les dataclasses et le carnet,
écrire les tests de conservation, ajouter les sept phases, exécuter tous les
tests, puis seulement lancer baseline et validations. Toute divergence entre
spécification et code devra être consignée dans le journal.

\section{Conclusion}

Cette spécification rend visibles les hypothèses logicielles capables de devenir
des résultats économiques. La fidélité sera évaluée par les invariants et non par
la ressemblance des graphes au prototype. La validation statistique sépare seeds,
fenêtres et variables, inclut des réfutations et interdit le réglage après coup.
Elle fournit ainsi une base falsifiable pour décider si M2 produit réellement un
corps exponentiel, une queue stable, des classes dynamiques et une population
bornée sans capacité artificielle.

\subsection*{Amendement après le premier run long}

Le run de référence a révélé une explosion purement représentationnelle des
fractions de contrats : 3\,482 arêtes à $t=1200$, 25\,243 à $t=1300$, 74\,807 à
$t=1400$ et 142\,766 à $t=1414$. Le carnet fusionne donc désormais tous les
contrats d'une même paire prêteur--emprunteur, avec principal total et taux moyen
pondéré. Cette compression conserve exactement principal, intérêt total futur,
créance, dette et poids prorata, puisque les contrats sont nominaux, perpétuels,
sans maturité ni dépréciation. Elle perd une partie de la lignée individuelle des
fragments, remplacée par un compteur de composantes. Il s'agit d'une correction
de représentation documentée, non d'une règle économique ni d'un seuil de
troncature.

\end{document}
